Notre \robot\ peut sauter au bord ou au centre du cadre (qui délimite
la zone de dessin sur l'écran de l'ordinateur), sans changer
son orientation, ni rien dessiner. C'est exactement comme si nous pouvions
 le porter
à la main pour le poser au bord ou au centre du cadre comme on peut le faire avec une voiture
télécommandée. L'ordinateur le fait à notre place grâce aux commandes 
données dans la table ci-dessous et illustrées à la figure
\ref{fig:ecrangdhbc} (page \pageref{fig:ecrangdhbc}) :
\begin{table}[htbp]
\begin{center}
\begin{tabular}{@{}|l|p{11cm}|@{}} \hline
{\tt ag()} & place le \robot\ \underline{à} \underline{g}auche de l'écran, sans changer sa
                  position verticale, ni son orientation, et sans rien dessiner.\\ \hline
{\tt ad()} & place le \robot\ \underline{à} \underline{d}roite de l'écran, sans changer sa
                  position verticale, ni son orientation, et sans rien dessiner.\\ \hline
{\tt eh()} & place le \robot\ \underline{e}n \underline{h}aut de l'écran, sans changer sa
                  position horizontale, ni son orientation, et sans rien dessiner.\\ \hline
{\tt eb()} & place le \robot\ \underline{e}n \underline{b}as de l'écran, sans changer sa
                  position horizontale, ni son orientation, et sans rien dessiner.\\ \hline
{\tt ce()} & place le \robot\ au \underline{ce}ntre de l'écran, sans changer son orientation, et
                  sans rien dessiner.\\ \hline
\end{tabular}
\end{center}
%\label{tab:ag-ad-eh-eb-ce}
\caption{Sauts du \protect\robot\ au bord et au centre du cadre}
\end{table}

{\samepage %1
\begin{figure}[ht]
\includevspacePostScript[7.2cm]{ecran-ag-ad-eh-eb-ce}{0.95}{0.95}
\caption{Sauts du \protect\robot\ au bord et au centre du cadre}
\label{fig:ecrangdhbc}
\end{figure}
\noindent \rule{137mm}{0.1mm} %---------------------------------------------
\begin{exercice}{
Écrire un programme \python\ qui dessine des croix au bord et au centre de l'écran
comme le montre  la figure ci-dessous.
\noindent $\Box$}\end{exercice}
\vspace*{15mm}
\includevspacePostScript[6.5cm]{CroixDeCoins}{0.95}{0.95}
\begin{figure}[hb]
 \caption{Croix au bord et au centre du cadre}
 %\label{fig:CroixDeCoins}
\end{figure}
} %1


\newpage
\begin{corrige}{\ \\

\vspace{-\baselineskip}
\input python/CroixDeCoins.tex
\noindent $\Box$}\end{corrige}

%%%%%%%%%%%%%%%%%%%%%%%%%%%%%%%%%%%%%%%%%%%%%%%%%%%%%%%%%%%%%%%%%%%%%%%%%%%
\noindent \framebox[\textwidth]{
\begin{minipage}{\myparboxwidth}
\setlength{\parindent}{2em}
\small \sf
\vspace{\abovedisplayskip}
\section*{Addenda}

Les commandes de déplacement du \robot\ au bord du cadre seront définies en \python\ ultérieurement après
avoir expliqué les notions de test et d'itération. Nous les introduisons dès maintenant pour pouvoir dessiner aisément 
des frises.  Dans le chapitre suivant ces frises à motifs répétitifs vont nous permettre d'introduire simplement
la notion d'itération bornée.

La commande {\idt ce()} de recentrage du \robot\ sera essentielle dès que nous utili\-se\-rons un repère cartésien pour désigner les points
de la grille.

\vspace{\abovedisplayskip}
\end{minipage}
}
%%%%%%%%%%%%%%%%%%%%%%%%%%%%%%%%%%%%%%%%%%%%%%%%%%%%%%%%%%%%%%%%%%%%%%%%%%%%
\newpage
\thispagestyle{empty}